\section{Polygraphic homology}
\begin{paragr}
  Recall that we have a so-called ``abelianization'' functor $\abel
  \colon \oCat \to \Chp$ from the category of $\oo$\nbd-categories to
  the category of chain complexes in non-negative degree, defined in
  the following way. For an $\oo$\nbd-category $C$ and $n\geq 0$
  \[
    \abel_n(C):=\Z C_n/\sim_n,
  \]
  where $\Z C_n$ is the free abelian group on the set of $n$\nbd-cells
  of $C$ and $\sim_n$ is the smallest equivalence relation compatible with
  the abelian group structure such that $x\comp{i} y
  \simeq _n x + y$ for any $x,y\in C_n$ that are
  $\comp{i}$-composables. The boundary operator is defined as
  \[
    \begin{aligned}
      \partial \colon \abel_{n+1}(C) &\to \abel_{n}(C)\\
      [x] &\mapsto [\target x] - [\source x].
    \end{aligned}
  \]
  One checks that this is well-defined, that we have $\partial
  \circ \partial=0$ and that the correspondence $C \mapsto
  \abel_{\bullet}(C)$ can be made functorial in an obvious
  way. Furthermore, $\abel$ admits a right adjoint. For
  details we refer to \cite{abgmmm:polybk}.

  By pre-composing with the inclusion functor, for any $k\geq 0$, we
  also get an abelinization functor
  \[
    \abelk{k} \colon \opCat{k} \overset{\emb{k}{\oo}}{\longrightarrow} \oCat \overset{\lambda}{\longrightarrow} \Chp.
  \]
  In particular, for $k=\oo$, we have $\abel=\abelk{\oo}$ by
  definition. Notice that since $\emb{k}{\oo}$ admits a right adjoint, so does $\abelk{k}$.
\end{paragr}

The abelianization of \ook k-polygraphs admits a useful explicit
description as stated in the following result.
\begin{proposition}
  Let $k \geq 0$, $P$ an \ook k-polygraph and $P_{n}$ a set of
  generators for each $n \geq 0$. Then $\abel_{\bullet}(\frgp P)$
  is canonically isomorphic to the chain complex
  \[
    \Z P_0 \overset{\partial}{\longleftarrow} \Z P_1
    \overset{\partial}{\longleftarrow} \dots,
  \]
  where $\partial$ is defined on generator by
  \[
    \partial[x] = [\target x] - [\source x].
  \]
\end{proposition}
\lgcomment{A-t-on besoin de cette description? Par ailleurs, la définition de l'opérateur bord dans cette proposition
  n'est pas très précise.}
\begin{proof}
\end{proof}
\begin{paragr}
  Let $0 \leq k
  \leq k' \leq \omega$. We are now going to analyze the behaviour of
  the abelianization functor with respect to the left adjoint $\pifun{k'}{k} \colon
  \opCat{k'} \to \opCat{k}$ of the inclusion functor. By definition,
  we have a commutative triangle
  \[
    \begin{tikzcd}[column sep=small]
      \opCat{k}\ar[rr,"\emb{k}{k'}"]\ar[rd,"\abelk{k}"']&& \opCat{k'}\ar[dl,"\abelk{k'}"] \\
      &\Chp&.
  \end{tikzcd}
  \]
  Now, consider the
  following, \emph{a priori} not commutative, triangle
  \[
    \begin{tikzcd}[column sep=small]
      \opCat{k'}\ar[rr,"\pifun{k'}{k}"]\ar[rd,"\abelk{k'}"']&& \opCat{k}\ar[dl,"\abelk{k}"] \\
    &\Chp&.
  \end{tikzcd}
\]
The unit defines a natural transformation
\[
 \abelk{k'}\Longrightarrow \abelk{k'} \emb{k}{k'}\pifun{k'}{k}= \abelk{k}\pifun{k'}{k}
\]
\end{paragr}
\begin{lemma}\label{lemma:abelpi}
  The above natural transformation $\abelk{k'} \Rightarrow \abelk{k}\pifun{k'}{k}
  $ is an isomorphism. 
\end{lemma}
\begin{proof}
\end{proof}
\begin{paragr}
  Recall that the category $\oCat$ admits a monoidal category
  structure with tensor product the so-called ``Gray tensor product'',
  simply denoted by $\gray$, and with unit $\Dsk{0}$ (see for example
  \cite{}). More generally, for any $0 \leq k \leq \omega$, we have
  the induced tensor product $\grayk{k}$
\end{paragr}
\begin{proposition}
  For any $0 \leq k \leq \omega$, the abelianization functor
  \[
    \abelk{k} \colon \opCat{k} \to \Chp
  \]
  is strong monoidal with respect to the Gray tensor product $\grayk{k}$ on
  $\opCat{k}$ and the usual tensor product of chain complexes on $\Chp$.
\end{proposition}
\begin{proof}
  The case $k=\omega$ is standard (see for example \cite{}). The
  general case follows from  $\abelk{k}\simeq
  \abelk{\omega}\pifun{\oo}{k}$ (\cref{lemma:abelpi}) and the fact
  that $\abelk{\omega}$ and $\pifun{\oo}{k}$ are strong monoidal.
\end{proof}
\begin{proposition}
  For any $k \geq 0$, the abelianization functor
  \[
    \abelk{k} \colon \opCat{k} \to \Chp
  \]
  is left Quillen, when $\npCat{\oo}{k}$ is equipped with the folk
  model structure and $\Chp$ with the projective model structure.
\end{proposition}
\begin{proof}
    The case $k=\omega$ is \cite[Theorem 22.2.7]{abgmmm:polybk}. For the
  general case, since the folk model structure on $\opCat{k}$ is right induced from
  the inclusion functor $\emb{k}{\omega} \colon \opCat{k} \to \oCat
  $, this means that a set of generating cofibrations (resp.\ trivial
  cofibrations) is given by the image by $\pifun{k}{\omega} \colon \oCat
  \to \opCat{k}$ of generating cofibrations of
  $\pifun{k}{\omega}$. The result follows then immediatly from \cref{lemma:abelpi}.
\end{proof}
\begin{paragr}
  The previous proposition contains an important subtlety. % For any $0
  % \leq k \leq k' \leq \omega$, the inclusion functor
  % $\embed{k}{k'}\colon \opCat{k} \to \opCat{k}$ is
  Even though for any  $0 \leq k \leq k' \leq \omega$ we have a
  commutative triangle,
  \[
    \begin{tikzcd}[column sep=small]
      \opCat{k}\ar[rr,"\emb{k}{k'}"]\ar[rd,"\abelk{k}"']&& \opCat{k'}\ar[dl,"\abelk{k'}"] \\
      &\Chp&,
  \end{tikzcd}
\]
and even though $\emb{k}{k'}$ passes to the homotopy categories, as
it trivially preserves the weak equivalences, this doesn't \emph{a
  priori} imply that the following triangle at the level of homotopy categories
\[
  \begin{tikzcd}[column sep=small]
    \Loc{\opCat{k}}\ar[rr,"\emb{k}{k'}"]\ar[rd,"\LL \abelk{k}"']&&
    \Loc{\opCat{k'}}\ar[dl,"\LL \abelk{k'}"] \\
    &\Loc{\Chp}&,
  \end{tikzcd}
\]
is also commutative, and the universal property of left derived functors
only yields a natural transformation
\begin{equation}\label{transnat}
  \LL \abelk{k'}\circ \emb{k}{k'} \Rightarrow \LL \abelk{k},
\end{equation}
which is \emph{a priori} not an isomorphism (and we do not know if it
is in general, see \cref{openquestion}.)

More concretly, this natural transformation can be defined as follows:
let $C$ be an \ook k\nbd-category and let
\[
  P\to C
\]
be a cofibrant resolution in $\opCat{k}$. Since $\emb{k}{k'}$ does
not preserve cofibrant objects in general, we need to take a cofibrant
resolution in $\opCat{k'}$
\[
  P' \to \emb{k}{k'}P.
\]
The above natural transformation is then
\[
  \LL \abelk{k'}(\emb{k}{k'}(C))=\abelk{k'}(P') \longrightarrow
  \abelk{k'}(\emb{k}{k'}(P))=\abelk{k}(P)=\LL \abelk{k}(C).
\]
\end{paragr}

\begin{definition}
  Let $k\geq 0$. The \emph{$\pair \oo k$\nbd-polygraphic homology} of
  an \ook k\nbd-category $C$ is defined as
  \[
    \Hop{k}_{\bullet}(C):=\Ho_{\bullet}(\LL \abelk{k} (C)).
  \]
\end{definition}

\begin{paragr}
  For $0 \leq k \leq k'\leq \omega$ and an \ook k\nbd-category $C$, we shall
  simply write $\Hop{k'}_{\bullet}(C)$ instead of
  $\Hop{k'}_{\bullet}(\emb{k}{k'}(C))$.
  In particular, the natural transformation \eqref{transnat} becomes
  \[
    \Hop{k'}_{\bullet}(C) \to \Hop{k}_{\bullet}(C).
  \]
\end{paragr}

\begin{openquestion}\label{openquestion} With the same notation as
  above, is the canonical morphism
  \[
     \Hop{k'}_{\bullet}(C) \to \Hop{k}_{\bullet}(C)
   \]
   an isomorphism of $\Loc{\Chp}$?
\end{openquestion}
As we shall see at the end of this section, we answer positively to
the above open question in a very particular subcase. 

\begin{theorem}
  Let $C$ be a $\pk 1 i$\nbd-category with $i \in \{0,1\}$. Then, for any $i
  \leq k \leq \oo$ there is an isomorphism
  \[
    \Hop{k}_{\bullet}(C) \to \Ho_{\bullet}(C)
  \]
  natural in $C$. Furthermore, this isomorphism can be chosen
  naturally in $k$, meaning that for $i \leq k \leq k' \leq \oo$, the
  triangle
  \[
    \begin{tikzcd}[column sep=small]
         \Hop{k'}_{\bullet}(C)\ar[dr]
         \ar[rr]&&\Hop{k}_{\bullet}(C)\ar[dl]\\
         &\Ho_{\bullet}(C)&
    \end{tikzcd}
  \]
  is commutative. 
\end{theorem}
\begin{proof}
  We begin by the first part of the statement. Because $k \geq i$, we
  can see $C$ as an $\ook k$\nbd-category, via the canonical
  inclusion $\npCat{1}{i} \to \opCat{k}$, which we omit to denote. Let $P \to C$ be a
  polygraphic resolution in $\opCat{k}$. By 
\end{proof}
%%% Local Variables:
%%% mode: LaTeX
%%% TeX-master: "main"
%%% End:
